% 2000 春季京皖理科第14题。示意性拓扑重绘，不按比例，不表达精确函数系数。
\begin{tikzpicture}
\begin{axis}[
  width=6.3cm,height=3.7cm,font=\small,
  axis lines=middle,xmin=-0.45,xmax=2.65,ymin=-1.25,ymax=1.15,
  xtick=\empty,ytick=\empty,
  xlabel={$x$},ylabel={$y$},
  every axis x label/.style={at={(current axis.right of origin)},anchor=west},
  every axis y label/.style={at={(current axis.above origin)},anchor=south}]
  % Cubic sketch only: visible crossings and axis labels are retained; scale is arbitrary.
  \draw (axis cs:0,-0.035) -- (axis cs:0,0.035);
  \draw (axis cs:1,-0.035) -- (axis cs:1,0.035);
  \draw (axis cs:2,-0.035) -- (axis cs:2,0.035);
  \node[below=2pt,inner sep=1pt] at (axis cs:0,0) {$0$};
  \node[below=2pt,inner sep=1pt] at (axis cs:1,0) {$1$};
  \node[below=2pt,inner sep=1pt] at (axis cs:2,0) {$2$};
  \addplot[domain=-0.3:2.3,samples=120,thick] {0.42*(x)*(x-1)*(x-2)};
\end{axis}
\end{tikzpicture}
